% 已完成的工作 — 每天更新进度时编辑此文件

\subsection{macOS 用户态基础}

\begin{itemize}
    \item 程序入口 \texttt{\_start} 与主循环骨架
    \item 条件跳转（\texttt{cmp}、\texttt{b.eq} 等）
    \item macOS exit 系统调用
    \item LLDB 调试：查看寄存器、设置断点、单步、检查内存
    \item 多文件拆分与 \texttt{.global/.extern} 跨文件调用
\end{itemize}

\subsection{内存管理}

\begin{itemize}
    \item 堆结构定义：\texttt{heap}、\texttt{heap\_end}、\texttt{HEAP\_SIZE}
    \item 块头格式：约定 size + used flag
    \item 分配扫描与分割：遍历堆块、标记已使用、大块分裂
    \item 释放与相邻合并：向前合并和向后合并
    \item 边界检查：空间不足时安全返回
\end{itemize}

\subsection{系统调用与 I/O}

\begin{itemize}
    \item write / read 系统调用封装
    \item 单字符输出与读取
    \item 字符串输出与读取
    \item 输入/输出缓冲区定义
\end{itemize}

\subsection{字符串运行时}

\begin{itemize}
    \item 堆字符串复制：从输入缓冲区拷贝到堆空间
    \item 平行指针/长度数组：\texttt{string\_ptrs} 存指针，
          \texttt{string\_lens} 存长度
    \item 按序号存取：通过 index 获取对应字符串的指针和长度
    \item 按序号打印：取回 ptr/len 后调用 write
\end{itemize}

\subsection{解析器前端}

\begin{itemize}
    \item \texttt{find\_two\_spaces}：逐字节扫描字符串，记录两个空格的位置
    \item 失败路径：空格不足时返回错误值
\end{itemize}
